\documentclass[spanish,letterpaper,oneside]{article}

\def\documenttitle {Actividad Diagnóstica: Enseñanza, Aprendizaje, Paradigmas y Matemáticas}
\def\documentsubtitle {Fundamentos para la enseñanza y el aprendizaje I · Subproducto 1A (Sesión 1)}
\def\documentsubject {Subproducto Independiente 1A · Diagnóstico de Cuatro Conceptos (Módulo 50157)}
\def\documentauthor {Mtro. Martín Jonathan de la Cruz Muñoz}
\def\coursename {Fundamentos para la enseñanza y el aprendizaje I}
\def\coursecode {MAEM-SD26}

\def\universityname {Instituto de Investigación, Innovación y Estudios de Posgrado para la Educación del Estado de Nuevo León}
\def\universityfaculty {Maestría en Aprendizaje y Enseñanza de las Matemáticas}
\def\universitydepartment {Fundamentos para la enseñanza y el aprendizaje I}
\def\universitydepartmentimage {img/iiiepe/marca-iiiepe.png}
\def\universitydepartmentimagecfg {height=1.75cm}
\def\universitylocation {Monterrey, Nuevo León}
\def\studentfullname {Mtro. Martín Jonathan de la Cruz Muñoz}
\def\studentid {00874}
\def\studentemail {martin.delacruz@campus.iiiepe.edu.mx}
\def\programperiod {2do periodo, septiembre--diciembre de 2026 (SD26)}
\def\programmodality {Modalidad presencial}
\def\programshort {MAEM}

\def\authortable {
  \begin{tabular}{ll}
    \textbf{Alumno:} & \studentfullname \\
    Matrícula: & \studentid \\
    Correo institucional: & \studentemail \\
    Unidad: & \coursename \\
    Programa: & \programshort \\
    Periodo: & \programperiod \\
    & \\
    \multicolumn{2}{l}{Fecha de realización: 19 de noviembre de 2026} \\
    \multicolumn{2}{l}{\universitylocation}
  \end{tabular}
}

\input{template}

\usepackage{helvet}
\renewcommand{\familydefault}{\sfdefault}
\usepackage{array}
\usepackage{longtable}
\usepackage{pdflscape}
\usepackage{tikz}
\usepackage{pgfplots}
\pgfplotsset{compat=1.18}
\usepackage[most]{tcolorbox}
\usetikzlibrary{arrows.meta,positioning,calc,fit,backgrounds,shapes.geometric}

\definecolor{iiiepeNavy}{HTML}{163B5C}
\definecolor{iiiepeBlue}{HTML}{1F6F9F}
\definecolor{iiiepeAqua}{HTML}{20B7A8}
\definecolor{iiiepeOrange}{HTML}{F28C28}
\definecolor{iiiepePaper}{HTML}{F6F8F7}
\definecolor{iiiepeInk}{HTML}{1E2A32}
\definecolor{iiiepePale}{HTML}{EEF5F8}
\definecolor{iiiepeMint}{HTML}{E7F6F4}
\definecolor{iiiepeWarm}{HTML}{FFF3E6}
\def\portraittitlecolor {iiiepeNavy}
\setcitestyle{authoryear,open={(},close={)}}
\makeatletter
\renewcommand{\@biblabel}[1]{}
\makeatother
\setlength{\bibhang}{1.27cm}
\setlength{\bibsep}{6pt}
\AtBeginDocument{\renewcommand{\refname}{Referencias}}

\newcommand{\insertcoverwatermark}{%
  \AddToShipoutPictureBG*{%
    \begin{tikzpicture}[remember picture,overlay]
      \node[opacity=0.04,inner sep=0pt] at (current page.center) {%
        \includegraphics[width=0.65\paperwidth]{img/iiiepe/marca-iiiepe.png}%
      };
    \end{tikzpicture}%
  }%
}

\setlength{\parskip}{5pt}\setlength{\parindent}{0pt}
\begin{document}\insertcoverwatermark\templatePortrait\templatePagecfg\templateFinalcfg\fancyfoot{}\fancyfoot[L]{\scriptsize IIIEPE · MAEM · SD26}\fancyfoot[C]{\scriptsize\thepage}\fancyfoot[R]{\scriptsize Subproducto 1A · Diagnóstico}\setstretch{1.3}\raggedbottom
\section{Introducción al diagnóstico conceptual}
Como punto de partida de la unidad de aprendizaje \textit{Fundamentos para la enseñanza y el aprendizaje I} (módulo \texttt{50157}, Sección 1), este documento presenta de manera independiente mi análisis diagnóstico sobre los cuatro conceptos estructurantes de nuestra formación en la Maestría en Aprendizaje y Enseñanza de las Matemáticas: \textbf{enseñanza, aprendizaje, paradigmas educativos y matemáticas}.

Mi postura es que explicitar las creencias previas heredadas de la formación técnica e ingenieril constituye una condición indispensable para transformar la docencia. En consecuencia, examino cómo transitar desde una racionalidad centrada en la transmisión mecánica de reglas hacia una mediación comprensiva, dialógica y socialmente justa dentro de CONALEP, CECyTE e ingeniería \citep{scott2015futuro,unesco2015replantear}.

\section{Articulación conceptual y contraste diagnóstico}
Desde mi trayectoria en mantenimiento industrial, mecatrónica e ingeniería, advierto que ninguno de estos cuatro conceptos opera de manera aislada: los paradigmas que asumimos —consciente o tácitamente— determinan qué valoramos como aprendizaje válido y qué formas de enseñanza organizamos en torno a las matemáticas. La Figura~\ref{fig:sub01a-nexo} ilustra estas relaciones dinámicas.

\begin{figure}[htbp]
\centering
\begin{tikzpicture}[
  box/.style={draw=#1,fill=#1!10,thick,rounded corners=4pt,text width=4.65cm,minimum height=2.25cm,align=left,inner sep=7pt,font=\small},
  arr/.style={->,very thick,>=Latex,iiiepeNavy}
]
  \node[box=iiiepeNavy] (ens) at (0,3.1) {\textbf{\color{iiiepeNavy}1. Enseñanza}\\Diseño intencional de mediaciones, preguntas y andamiajes que habilitan el acceso equitativo al sentido matemático.};
  \node[box=iiiepeBlue] (apr) at (7.4,3.1) {\textbf{\color{iiiepeBlue}2. Aprendizaje}\\Construcción activa, situada y dialógica que integra comprensión conceptual, uso estratégico y transferencia.};
  \node[box=iiiepeAqua!80!black] (par) at (0,0) {\textbf{\color{iiiepeAqua!75!black}3. Paradigmas educativos}\\Marcos epistemológicos que determinan qué cuenta como saber válido, cómo se concibe el error y cómo se evalúa.};
  \node[box=iiiepeOrange!90!black] (mat) at (7.4,0) {\textbf{\color{iiiepeOrange!85!black}4. Matemáticas}\\Lenguaje de patrones, variación e inferencia para modelar fenómenos, tomar decisiones y argumentar con evidencia.};

  \draw[arr] (ens) -- node[above,font=\scriptsize\bfseries] {provoca y acompaña} (apr);
  \draw[arr] (par) -- node[left,font=\scriptsize\bfseries,align=right] {redefine\\metas} (ens);
  \draw[arr] (mat) -- node[right,font=\scriptsize\bfseries,align=left] {desafía\\creencias} (apr);
  \draw[arr] (par) -- node[below,font=\scriptsize\bfseries] {orienta el diseño} (mat);
\end{tikzpicture}
\caption{Diagnóstico relacional entre enseñanza, aprendizaje, paradigmas y matemáticas.}
\label{fig:sub01a-nexo}
\end{figure}

El cuadro comparativo siguiente contrasta la visión transmisiva inicial con la perspectiva comprensiva y humanista que asumo consolidar a lo largo de la cohorte SD26:

{\small\setstretch{1.15}\renewcommand{\arraystretch}{1.35}
\begin{longtable}{|>{\raggedright\arraybackslash}p{0.19\linewidth}|>{\raggedright\arraybackslash}p{0.37\linewidth}|>{\raggedright\arraybackslash}p{0.38\linewidth}|}\hline
\rowcolor{iiiepePale} \textbf{Concepto} & \textbf{Visión transmisiva o procedimental} & \textbf{Postura comprensiva y humanista} \\\hline
\endfirsthead
\hline
\rowcolor{iiiepePale} \textbf{Concepto} & \textbf{Visión transmisiva o procedimental} & \textbf{Postura comprensiva y humanista} \\\hline
\endhead
Enseñanza & Exposición unidireccional de reglas, fórmulas y ejemplos modelo seguidos de mecanización. & Mediación cognitiva que parte de problemas reales, estimula la conjetura y organiza el diálogo argumentado \citep{sawyer2008optimizing}. \\\hline
Aprendizaje & Retención de corto plazo y reproducción fiel de despejes sin comprensión del fenómeno. & Apropiación relacional y metacognitiva capaz de explicar significados y transferir métodos \citep{scott2015futuro}. \\\hline
Paradigmas & Racionalidad eficiencista centrada en cobertura temática y sanción punitiva del error. & Racionalidad socioconstructivista y humanista donde el error orienta el andamiaje y resguarda la equidad \citep{unesco2015replantear}. \\\hline
Matemáticas & Colección aislada de símbolos algebraicos desconectados de magnitudes y contexto. & Lenguaje vivo de modelación, argumentación y lectura crítica en entornos técnicos y sociales \citep{sep2023pensamiento}. \\\hline
\end{longtable}}

\section{Ejemplos matemáticos de lectura relacional en mi contexto docente}
Además, para volver observable este cambio conceptual en la formación tecnológica e ingeniería, identifico dos casos emblemáticos donde comprender supera la mera ejecución algorítmica:
\begin{enumerate}
  \item \textbf{Desplazamiento lineal con referencia inicial ($d(t) = 0.65\,t + 1.4$):} Frente a un carro de banco con sensor a $1.4\text{ m}$ de origen y velocidad constante $0.65\text{ m/s}$, calcular $d(8)=6.6\text{ m}$ es insuficiente si no se explica por qué la presencia de $d_0 \neq 0$ convierte el modelo en relación afín no directamente proporcional y se discute la longitud física del riel \citep{conalep2025mecatronica}.
  \item \textbf{Comparación de costos de mantenimiento ($C_A(h)=250+80h$ frente a $C_B(h)=100+110h$):} Determinar algebraicamente el punto de equilibrio $h=5\text{ h}$ adquiere sentido formativo cuando el estudiantado argumenta qué esquema conviene antes y después de cinco horas articulando tabla, gráfica, unidades ($\text{\$}/\text{h}$) y toma de decisiones \citep{star2015algebra}.
\end{enumerate}

\clearpage
\section{Conclusiones}
En conclusión, considero que este diagnóstico inicial aclara tanto el valor de la disciplina lógica técnica como la urgencia de complementarla con empatía pedagógica, problemas abiertos y seguimiento formativo. Sin embargo, dado que transformar hábitos expositivos arraigados no ocurre de modo automático, además resulta indispensable documentar qué obstáculo enfrenta cada grupo; es decir, por tanto sostengo que revisar continuamente qué entendemos por enseñar y aprender matemáticas constituye el primer acto ético para garantizar trayectorias educativas dignas y comprensivas \citep{delors1996educacion,unesco2015replantear}.

\clearpage
\begin{thebibliography}{99}
\bibitem[CONALEP(2025)]{conalep2025mecatronica} Colegio Nacional de Educación Profesional Técnica. (2025). \textit{Perfil de egreso: Profesional Técnico-Bachiller en Mecatrónica. Modelo Académico 2025}. Secretaría Académica, CONALEP.
\bibitem[Delors et~al.(1996)]{delors1996educacion} Delors, J., et al. (1996). \textit{La educación encierra un tesoro: Informe a la UNESCO de la Comisión Internacional sobre la Educación para el Siglo XXI}. Santillana/UNESCO.
\bibitem[Sawyer(2008)]{sawyer2008optimizing} Sawyer, R. K. (2008). Optimizing learning: Implications of learning sciences research. En OECD, \textit{Innovating to learn, learning to innovate} (pp. 45--65). OECD Publishing.
\bibitem[Scott(2015)]{scott2015futuro} Scott, C. L. (2015). \textit{El futuro del aprendizaje 1: ¿Por qué deben cambiar el contenido y los métodos de aprendizaje en el siglo XXI?} UNESCO.
\bibitem[SEP(2023)]{sep2023pensamiento} Secretaría de Educación Pública. (2023). \textit{Programa de estudios del recurso sociocognitivo Pensamiento Matemático I}. SEP.
\bibitem[Star et~al.(2015)]{star2015algebra} Star, J. R., et al. (2015). \textit{Teaching strategies for improving algebra knowledge in middle and high school students}. IES.
\bibitem[UNESCO(2015)]{unesco2015replantear} UNESCO. (2015). \textit{Replantear la educación: ¿Hacia un bien común mundial?} Ediciones UNESCO.
\end{thebibliography}

\end{document}
